\documentclass[11pt]{article}
\usepackage[margin=1in]{geometry}
\usepackage{booktabs}
\usepackage{siunitx}
\usepackage{listings}
\usepackage{xcolor}
\usepackage{caption}

\lstset{
    language=C,
    basicstyle=\ttfamily\small,
    keywordstyle=\color{blue},
    commentstyle=\color{gray},
    numbers=left,
    numberstyle=\tiny\color{gray},
    frame=single,
    breaklines=true
}

\title{Performance Analysis: Complex Float vs Parametric Form\\for Mandelbrot Set Computation}
\author{}
\date{}

\begin{document}
\maketitle

\section{Introduction}

This analysis compares two implementations of the Mandelbrot set escape-time algorithm:
\begin{itemize}
    \item \textbf{serial\_0}: Uses C99 \texttt{float complex} type with \texttt{cabsf()} for magnitude calculation
    \item \textbf{serial\_1}: Uses explicit parametric form with separate \texttt{float} variables for real and imaginary components
\end{itemize}

Both implementations compute a $4000 \times 4000$ grid with a maximum of 1000 iterations per point.

\section{Implementation Comparison}

\subsection{Complex Float Implementation (serial\_0)}

\begin{lstlisting}[caption={Hot loop using complex float},firstnumber=28]
k = 1;
while ((cabsf(z) <= 2.0f) && (k++ < MAXITER)) 
    z = z * z + kappa;
\end{lstlisting}

\subsection{Parametric Form Implementation (serial\_1)}

\begin{lstlisting}[caption={Hot loop using parametric form},firstnumber=29]
k = 1;
while (((zi * zi) + (zj * zj) <= 4.0f) && (k++ < MAXITER)) { 
    new_zi = (zi * zi) - (zj * zj) + ki;
    zj = 2.0f * zi * zj + kj;
    zi = new_zi;
}
\end{lstlisting}

\section{Profiling Methodology}

Programs were compiled with \texttt{gcc -O3 -g -pg} and profiled using:
\begin{itemize}
    \item \textbf{gprof} with line-level annotation (\texttt{-l} flag) to attribute execution time to individual source lines
    \item \textbf{perf stat} to measure hardware performance counters (cycles, instructions, cache misses)
\end{itemize}

\section{Results}

\subsection{gprof Line-Level Analysis}

Table~\ref{tab:gprof} shows the percentage of execution time attributed to each source line in the hot loop.

\begin{table}[h]
\centering
\caption{gprof line-level profiling results}
\label{tab:gprof}
\begin{tabular}{@{}lrrr@{}}
\toprule
\textbf{Source Line} & \textbf{serial\_0} & \textbf{serial\_1} & \textbf{serial\_2} \\
\midrule
Line 29: while condition (\texttt{cabsf}) & 87.83\% & --- & --- \\
Line 30: \texttt{z = z*z + kappa} / while condition & 5.83\% & 73.58\% & 69.99\% \\
Line 31: \texttt{new\_zi} calculation & --- & 3.11\% & 6.18\% \\
Line 32: \texttt{zj} update & 0.16\% & 13.62\% & 15.45\% \\
Line 36: \texttt{logf()} normalization & --- & 9.32\% & 7.64\% \\
\midrule
\textbf{Total self time} & \textbf{12.43s} & \textbf{5.47s} & \textbf{2.75s} \\
\bottomrule
\end{tabular}
\end{table}

\subsection{Hardware Performance Counters}

Table~\ref{tab:perf} presents the hardware-level metrics collected via \texttt{perf stat}.

\begin{table}[h]
\centering
\caption{perf stat hardware performance metrics}
\label{tab:perf}
\sisetup{group-separator={,}, group-minimum-digits=4}
\begin{tabular}{@{}lSSS@{}}
\toprule
\textbf{Metric} & \textbf{serial\_0} & \textbf{serial\_1} & \textbf{serial\_2} \\
\midrule
Cycles (billions) & 94.9 & 13.6 & 6.9 \\
Instructions (billions) & 89.9 & 28.4 & 14.3 \\
IPC (instructions/cycle) & 0.95 & 2.08 & 2.08 \\
Wall time (seconds) & 38.2 & 5.5 & 2.8 \\
\midrule
\textbf{Speedup vs serial\_0} & {1.0x} & {6.9x} & {13.6x} \\
\bottomrule
\end{tabular}
\end{table}

\section{Analysis}

\subsection{Why Complex Float is Slow}

The gprof results reveal that \textbf{87.8\% of serial\_0's execution time} is spent on line 29, which contains the \texttt{cabsf(z)} call. This function computes:
\[
|z| = \sqrt{z_r^2 + z_i^2}
\]

The \texttt{sqrtf()} operation is computationally expensive, requiring either iterative approximation or specialized hardware instructions with high latency. Furthermore, the comparison \texttt{cabsf(z) <= 2.0f} could equivalently be written as $z_r^2 + z_i^2 \leq 4.0$, avoiding the square root entirely.

The low IPC of \textbf{0.95} indicates that the CPU pipeline is frequently stalled waiting for the \texttt{sqrtf()} result, preventing effective instruction-level parallelism.

\subsection{Why Parametric Form is Faster}

The parametric implementation achieves a \textbf{6.9x speedup} through two key optimizations:

\begin{enumerate}
    \item \textbf{Elimination of \texttt{sqrtf()}}: The escape condition \texttt{zi*zi + zj*zj <= 4.0f} avoids the expensive square root by comparing squared magnitudes directly.
    
    \item \textbf{Explicit arithmetic enables better optimization}: The compiler can more effectively schedule the simple floating-point multiply-add operations, achieving an IPC of \textbf{2.08}---more than double that of the complex version.
\end{enumerate}

The total instruction count drops from 89.9B to 28.4B (\textbf{3.2x reduction}), and each instruction executes faster due to improved pipelining.

\subsection{Symmetry Optimization (serial\_2)}

The serial\_2 implementation adds an additional \textbf{2x speedup} by exploiting the Mandelbrot set's symmetry about the real axis. Since $f(z) = f(\bar{z})$, only half the points need to be computed. This halves both the instruction count and execution time while maintaining identical IPC.

\section{Conclusion}

The parametric form implementation outperforms the complex float version by \textbf{6.9x} due to:
\begin{itemize}
    \item Avoiding the expensive \texttt{sqrtf()} call hidden inside \texttt{cabsf()}
    \item Enabling better instruction-level parallelism (IPC: 2.08 vs 0.95)
    \item Reducing total instructions by 3.2x
\end{itemize}

For performance-critical numerical code, explicit arithmetic on primitive types often outperforms abstracted types like \texttt{complex}, even when the latter produces cleaner source code.

\end{document}
